% ==========================================================================
% GENERATED by scripts/split_paper.py. Do not edit this file.
%
% Source: imported_papers/grokking/iclr2027_conference_v2.tex
% Part 12: Other tasks and architectures
%
% Edits belong upstream, in the published paper. To change what the
% chapter includes, change imports/grokking.toml and re-run
% the splitter.
% ==========================================================================
\section{Other tasks and architectures}
\label{app:scope}
The main experiments use modular addition and a multilayer perceptron. Here we repeat the comparison between the reference and full credit equalization on other operations, tasks and architectures (Table~\ref{tab:scope}). The intervention raises participation where the reference circuit is concentrated, with one exception, layer normalization, and it changes little where the reference circuit is already distributed.

\begin{table}[!ht]
\caption{Circuit participation of the reference ($\beta{=}0$) and under full credit equalization ($\beta{=}1$) across the tested settings. Test accuracies and seed counts are given in the cited subsections. Sparse parity and MNIST show no grokking delay. The last three settings are unchanged: two start from a distributed circuit, and with layer normalization the reference is concentrated.}
\label{tab:scope}
\begin{center}
\footnotesize
\setlength{\tabcolsep}{5pt}
\begin{tabular}{llccl}
\toprule
Setting & Circuit measured & $\beta{=}0$ & $\beta{=}1$ & Details \\
\midrule
Addition, subtraction, multiplication & first hidden layer & $0.43$--$0.60$ & $0.99$ & App.~\ref{app:ops} \\
Sparse parity           & first hidden layer & $0.15$ & $0.96$ & App.~\ref{app:parity} \\
$S_5$ composition       & first hidden layer & $0.70$ & $0.97$ & App.~\ref{app:groups} \\
Transformer, $p=31$, wide MLP & MLP block & $0.61$ & $0.99$ & App.~\ref{app:transformer} \\
Batch normalization & first hidden layer & $0.81$ & $0.97$ & App.~\ref{app:normlayers} \\
\midrule
Transformer, $p=113$, canonical & MLP block & $0.96$ & $0.96$ & App.~\ref{app:transformer} \\
MNIST                   & first hidden layer & $0.99$ & $0.99$ & App.~\ref{app:mnist} \\
Layer normalization & first hidden layer & $0.36$ & $0.37$ & App.~\ref{app:normlayers} \\
\bottomrule
\end{tabular}
\end{center}
\end{table}

\subsection{Subtraction and multiplication}
\label{app:ops}
Table~\ref{tab:ops} repeats the comparison on subtraction and multiplication at $p = 31$, with DFC for comparison. Equalization raises participation to $0.99$ on both operations while test accuracy stays at or above $0.994$, and random-deletion tolerance increases as on addition. The largest change is on multiplication, where the reference circuit is the most concentrated of the three, with participation $0.43$.

\begin{table}[!ht]
\caption{Circuit participation and random-deletion tolerance on the three modular operations ($p=31$, three seeds; standard deviations in parentheses). All three operations use the same random split with $40\%$ of the pairs held out.}
\label{tab:ops}
\begin{center}
\small
\begin{tabular}{llccc}
\toprule
Operation & Rule & Test & Circuit PR & Random del \\
\midrule
Addition       & reference & $0.998\,(0.001)$ & $0.60\,(0.02)$ & $0.671\,(0.016)$ \\
               & equalization ($\beta{=}1$) & $0.976\,(0.011)$ & $0.99\,(0.00)$ & $0.771\,(0.019)$ \\
               & DFC & $1.000\,(0.000)$ & $0.92\,(0.00)$ & $0.717\,(0.014)$ \\
\midrule
Subtraction    & reference & $0.999\,(0.001)$ & $0.59\,(0.04)$ & $0.680\,(0.026)$ \\
               & equalization ($\beta{=}1$) & $0.994\,(0.005)$ & $0.99\,(0.00)$ & $0.785\,(0.014)$ \\
               & DFC & $1.000\,(0.000)$ & $0.93\,(0.01)$ & $0.735\,(0.010)$ \\
\midrule
Multiplication & reference & $0.999\,(0.001)$ & $0.43\,(0.07)$ & $0.666\,(0.035)$ \\
               & equalization ($\beta{=}1$) & $0.999\,(0.001)$ & $0.99\,(0.00)$ & $0.867\,(0.012)$ \\
               & DFC & $1.000\,(0.000)$ & $0.86\,(0.02)$ & $0.764\,(0.010)$ \\
\bottomrule
\end{tabular}
\end{center}
\end{table}

\subsection{Network width}
\label{app:arch}
Does the difference between the two circuits depend on how many neurons are available? We vary the hidden width from $32$ to $2048$ with the task, the budget and the comparison between $\beta = 0$ and $\beta = 1$ fixed (Table~\ref{tab:widths}). At width $32$ the effective counts are $30$ for the reference and $31$ under equalization. At width $2048$ they are $332$ and $2046$. The reference uses a decreasing fraction of the layer as the width grows, from participation $0.937$ to $0.162$, while equalization stays near one. Test accuracy is at least $0.96$ throughout. Normalized participation and effective count answer different questions and can move in opposite directions: from width $128$ to $512$ the reference's participation falls from $0.791$ to $0.400$ while its effective count grows from $101$ to $205$. A wider reference network uses more neurons in absolute terms and a smaller share of the layer.

\begin{table}[!ht]
\caption{Width dependence of first-hidden-layer circuit participation and effective neuron count $n \cdot \mathrm{PR}$ at $\beta = 0$ and $\beta = 1$ (modular addition, $p=31$, three seeds, $15{,}000$ epochs).}
\label{tab:widths}
\begin{center}
\small
\begin{tabular}{rcccc}
\toprule
& \multicolumn{2}{c}{Circuit PR} & \multicolumn{2}{c}{$n \cdot \mathrm{PR}$} \\
Width & $\beta=0$ & $\beta=1$ & $\beta=0$ & $\beta=1$ \\
\midrule
$32$   & $0.937$ & $0.980$ & $30$  & $31$   \\
$64$   & $0.915$ & $0.987$ & $59$  & $63$   \\
$128$  & $0.791$ & $0.991$ & $101$ & $127$  \\
$512$  & $0.400$ & $0.994$ & $205$ & $509$  \\
$1024$ & $0.262$ & $0.994$ & $268$ & $1018$ \\
$2048$ & $0.162$ & $0.999$ & $332$ & $2046$ \\
\bottomrule
\end{tabular}
\end{center}
\end{table}

\subsection{Sparse parity}
\label{app:parity}
We test credit equalization on sparse parity, a task used to study how networks find a few relevant features after a long plateau \citep{barak2022hidden}. The input is a vector of $20$ bits, each $0$ or $1$; the label is $1$ if an odd number of three fixed bits are set and $0$ otherwise, so the remaining $17$ bits are distractors that the network has to ignore. We draw $4000$ or $600$ inputs at random from the $2^{20}$ possibilities and assign each to the training or the test set with probability one half, which gives about $2000$ and $300$ training examples. The runs generalize without delay under full-batch training, so this experiment tests the redistribution outside the grokking regime. Equalization raises circuit participation from $0.15$ to $0.96$ on the larger dataset and from $0.08$ to $0.89$ on the smaller one. Test accuracy is $1.000$ for both conditions on the larger dataset and changes from $1.000$ to $0.988$ on the smaller one, and random-deletion tolerance improves at both sizes (Table~\ref{tab:parity}). The redistribution is thus not restricted to modular arithmetic or to delayed generalization.

\begin{table}[!ht]
\caption{Sparse parity with approximately $2000$ and $300$ training examples (three seeds; standard deviations in parentheses).}
\label{tab:parity}
\begin{center}
\small
\begin{tabular}{llccc}
\toprule
Train size & Rule & Test & Circuit PR & Random del \\
\midrule
$\sim 2000$ & reference & $1.000\,(0.000)$ & $0.15\,(0.06)$ & $0.865\,(0.051)$ \\
            & equalization ($\beta{=}1$) & $1.000\,(0.000)$ & $0.96\,(0.03)$ & $0.955\,(0.002)$ \\
\midrule
$\sim 300$  & reference & $1.000\,(0.000)$ & $0.08\,(0.02)$ & $0.825\,(0.017)$ \\
            & equalization ($\beta{=}1$) & $0.988\,(0.017)$ & $0.89\,(0.02)$ & $0.908\,(0.060)$ \\
\bottomrule
\end{tabular}
\end{center}
\end{table}

\subsection{Permutation composition}
\label{app:groups}
We test whether credit equalization redistributes circuits on a noncommutative operation, the composition of permutations, which \citet{chughtai2023toy} use to study the circuits of grokked networks. The group $S_5$ contains the $120$ permutations of five objects. The input is a pair $(g, h)$ of permutations, and the target is their composition $g \circ h$, the permutation obtained by applying $h$ first and then $g$. We use one-hot inputs and train for $30{,}000$ epochs on half of the $120 \times 120$ input pairs, with the rest held out. Both the reference and the equalized network grok under the same definition as for modular arithmetic; the delay grows when fewer training pairs are available. Full equalization raises circuit participation from $0.699\,(0.078)$ to $0.973\,(0.005)$ while both conditions reach test accuracy $1.000$. Random-deletion tolerance increases from $0.485$ to $0.633$ and norm-ranked tolerance from $0.149$ to $0.468$ (Table~\ref{tab:groups}). Reducing the training fraction to $40\%$ reproduces the pattern. Credit and circuit participation again agree closely.

\begin{table}[!ht]
\caption{Composition in $S_5$. The upper block uses ten seeds with half the pairs held out, the lower block five seeds with $60\%$ held out; parentheses show standard deviations. Delay is the number of epochs between reaching training accuracy $1.0$ and test accuracy $0.95$.}
\label{tab:groups}
\begin{center}
\small
\begin{tabular}{lcccccc}
\toprule
Rule & Test & Credit PR & Circuit PR & Random del & Norm-ranked del & Delay \\
\midrule
\multicolumn{7}{l}{\emph{Half the pairs held out, ten seeds}} \\
Reference & $1.000$ & $0.689$ & $0.699\,(0.078)$ & $0.485$ & $0.149$ & $500$ \\
Equalization, $\beta{=}0.5$ & $1.000$ & $0.970$ & $0.942\,(0.010)$ & $0.585$ & $0.322$ & $550$ \\
Equalization, $\beta{=}1$ & $1.000$ & $1.000$ & $0.973\,(0.005)$ & $0.633$ & $0.468$ & $900$ \\
\midrule
\multicolumn{7}{l}{\emph{$60\%$ held out, five seeds}} \\
Reference & $1.000$ & $0.657$ & $0.662\,(0.038)$ & $0.409$ & $0.129$ & $1350$ \\
Equalization, $\beta{=}0.5$ & $1.000$ & $0.951$ & $0.907\,(0.016)$ & $0.525$ & $0.260$ & $1650$ \\
Equalization, $\beta{=}1$ & $1.000$ & $1.000$ & $0.970\,(0.002)$ & $0.597$ & $0.410$ & $1750$ \\
\bottomrule
\end{tabular}
\end{center}
\end{table}

\subsection{Transformers}
\label{app:transformer}
We test whether credit equalization changes the participation of a transformer's MLP block. We follow the modular-addition transformer of \citet{nanda2023progress}: inputs have the form ``$a\ b\ =$'', and the canonical configuration has one block, $p = 113$, $d_{\text{model}} = 128$, one attention head per $32$ channels, a ReLU MLP of width $4 d_{\text{model}}$ and weight decay $1.0$, trained with AdamW ($\beta_1 = 0.9$, $\beta_2 = 0.98$) at learning rate $10^{-3}$ on $30\%$ of the operand pairs for $40{,}000$ epochs. Credit equalization is applied to every linear map.

In the canonical model the reference MLP is already distributed: its participation is $0.955$, and full equalization leaves it at $0.955$ with test accuracy $1.000$ in both conditions (Table~\ref{tab:transformer}). A second modulus and a doubled model width give the same picture, with MLP participation changing by at most $0.015$. Applying the intervention only to the MLP changes participation by less than the variation across seeds. The intervention thus does not increase MLP participation where the reference is already distributed.

A wider MLP gives a concentrated reference. With $4096$ MLP neurons at $p = 31$, reference participation is $0.610\,(0.048)$, and full equalization raises it to $0.993\,(0.003)$ over five seeds, with the seed ranges not overlapping. Norm-ranked deletion tolerance rises from $0.580$ to $0.941$, test accuracy stays at $1.000$ and the margin is unchanged to two decimal places.

In the wide configuration, equalization acts on the MLP's incoming weights only, training uses $60\%$ of the operand pairs for $60{,}000$ epochs with a warmup followed by cosine decay, and AdamW decays every parameter, including the layer-normalization gains and shifts, as in the canonical setting. After about $20{,}000$ epochs the square roots of Adam's second-moment estimates for the MLP's incoming weights fall below $\epsilon = 10^{-8}$, to $10^{-11}$ to $5 \cdot 10^{-11}$ at the end, so Adam's step is approximately $\hat{\vm}/\epsilon$ and $\epsilon$ acts as the non-gradient term of the update. On five fresh seeds, the reference has participation $0.609$ and full equalization $0.994$. Dividing $\epsilon$ by each neuron's equalization factor, with the gradient unchanged, gives $0.981$. Equalizing the gradient and multiplying $\epsilon$ by the same factor, which keeps each neuron's ratio of gradient to $\epsilon$, gives $0.818$, so this control removes about half of the effect. With decay only on weight matrices, as in most current training code, $\epsilon$ no longer dominates. The reference is then broadly distributed, with participation $0.82$ over ten seeds, and equalization raises it to $0.91$; one of the ten equalized runs did not generalize.

Credit equalization thus has a large effect in the transformer whose reference MLP is concentrated and a small one where the reference is already distributed.

\begin{table}[!ht]
\caption{Modular-addition transformer results: test accuracy and MLP circuit participation (standard deviations in parentheses). Credit equalization is applied to every linear map; $\beta{=}0$ is the reference. The wide-MLP result is described in the text.}
\label{tab:transformer}
\begin{center}
\small
\begin{tabular}{lcc}
\toprule
$\beta$ & Test & MLP circuit PR \\
\midrule
\multicolumn{3}{l}{\emph{Canonical: one block, $p=113$, $d_{\text{model}}=128$, six seeds}} \\
$0.0$ & $1.000$ & $0.955\,(0.01)$ \\
$0.5$ & $1.000$ & $0.953\,(0.01)$ \\
$1.0$ & $1.000$ & $0.955\,(0.01)$ \\
\midrule
\multicolumn{3}{l}{\emph{Second modulus: $p=97$, three seeds, $70{,}000$ epochs}} \\
$0.0$ & $1.000$ & $0.959\,(0.01)$ \\
$1.0$ & $1.000$ & $0.951\,(0.02)$ \\
\midrule
\multicolumn{3}{l}{\emph{Twice the width: $d_{\text{model}}=256$, three seeds}} \\
$0.0$ & $1.000$ & $0.902\,(0.04)$ \\
$1.0$ & $1.000$ & $0.917\,(0.01)$ \\
\bottomrule
\end{tabular}
\end{center}
\end{table}

\subsection{MNIST}
\label{app:mnist}
We use the setting in which image-classification grokking was reported \citep{liu2022omnigrok}: $1000$ training images and a network with two hidden layers of width $200$, at the standard and at a larger initialization scale. Training and test accuracy rise together, so these runs show no grokking delay. The reference already has circuit participation $0.992\,(0.002)$, and full credit equalization changes it by less than the variation across seeds. As in the canonical transformer, the reference leaves little room for a further increase in participation.

\subsection{Normalization layers}
\label{app:normlayers}
We insert batch normalization \citep{ioffe2015batchnorm} or layer normalization \citep{ba2016layernorm} after each hidden linear map of the multilayer perceptron, before the ReLU, and otherwise train as in Section~\ref{sec:setup} for $20{,}000$ epochs over five seeds. Measurements use evaluation mode, with the running statistics for batch normalization. Batch normalization, which normalizes each neuron's pre-activation separately, partly distributes the circuit, to participation $0.808$ against $0.611$ without normalization, with $42$ instead of $81$ nearly unused neurons, and equalization raises it to $0.970$ (Table~\ref{tab:normlayers}). Layer normalization, which normalizes across the neurons of a layer, concentrates the circuit further, to participation $0.356$ with $164$ nearly unused neurons, and equalizing the first layer's gradients leaves it at $0.369$; on the five fresh seeds of Table~\ref{tab:pathways} the two values are $0.366$ and $0.364$. Test accuracy is $1.000$ with either normalization layer. Layer normalization concentrates the circuit through its trainable per-neuron gain, not through the normalization. With the normalization layer on the first hidden layer only, a trainable gain and bias without any normalization give the same participation as layer normalization, $0.099$ against $0.094$ with $231$ nearly unused neurons each, and RMS normalization gives $0.094$. Layer normalization without a trainable gain gives $0.53$ to $0.58$, and batch normalization without affine parameters gives $0.996$ with no intervention. Under layer normalization a neuron has two decayed scales of its own, the gain and the outgoing weights, and equalizing the incoming gradients reaches neither. Rescaling the decay of both by the inverse of each neuron's equalization factor distributes the circuit in five fresh seeds, with participation $0.997$ and no nearly unused neurons when layer normalization follows both hidden maps, and $1.000$ at test accuracy $0.94$ to $1.00$ when it follows the first. Equalizing the gradients of the incoming rows, the gain and the outgoing weights does the same, with $0.996$ and $1.000$, while equalizing them with each neuron's ratio of gradient to decay held fixed leaves participation at $0.343$ and $0.129$. In an earlier set of five seeds the activation-weighted participation follows, at $0.95$ and $0.98$ with the decay rescaled and $0.31$ and $0.12$ with the ratio held. With layer normalization on the first hidden layer, the circuit keeps about $26$ neurons at widths $64$, $256$ and $1024$, against $61$, $175$ and $305$ without normalization.

\begin{table}[!ht]
\caption{Normalization layers after each hidden linear map (modular addition, $p=31$, five seeds; standard deviations in parentheses). Columns as in Table~\ref{tab:decay}.}
\label{tab:normlayers}
\begin{center}
\footnotesize
\setlength{\tabcolsep}{4pt}
\begin{tabular}{lccccccc}
\toprule
Normalization, $\beta$ & Test & Circuit PR & Nearly unused & Act. & Abl. & Random & Norm-ranked \\
\midrule
None, $0$ & $0.993$ & $0.611\,(0.015)$ & $81.0$ & $0.338$ & $0.417$ & $0.669$ & $0.292$ \\
None, $1$ & $0.974$ & $0.991\,(0.003)$ & $0$ & $0.978$ & $0.916$ & $0.776$ & $0.594$ \\
Batch, $0$ & $1.000$ & $0.808\,(0.007)$ & $42.4$ & $0.735$ & $0.698$ & $0.805$ & $0.556$ \\
Batch, $1$ & $1.000$ & $0.970\,(0.003)$ & $0.2$ & $0.892$ & $0.845$ & $0.834$ & $0.651$ \\
Layer, $0$ & $1.000$ & $0.356\,(0.004)$ & $164.0$ & $0.336$ & $0.307$ & $0.584$ & $0.195$ \\
Layer, $1$ & $1.000$ & $0.369\,(0.005)$ & $161.0$ & $0.352$ & $0.328$ & $0.605$ & $0.214$ \\
\bottomrule
\end{tabular}
\end{center}
\end{table}

\subsection{A language model}
\label{app:lm}
We test whether the per-neuron condition decides which MLP units survive in a language model. The model is a GPT-2-style transformer with layer normalization before each block, six blocks, model width $384$, MLP width $1536$ with GELU, context $256$ and tied embeddings, about $30$M parameters. It is trained on TinyStories \citep{eldan2023tinystories} for $393$M tokens ($24{,}000$ steps of $64$ sequences) with Adam ($\beta_1 = 0.9$, $\beta_2 = 0.95$), learning rate $2 \cdot 10^{-3}$ with warmup and cosine decay, gradient clipping at $1.0$ and coupled L2 decay $3 \cdot 10^{-5}$ on all weight matrices. As in the multilayer perceptrons, a unit is nearly unused if its outgoing norm is below a tenth of its block's mean. The balanced condition scales the decay of each unit's incoming row and outgoing column by the ratio of that unit's gradient norm to the block mean, which leaves the mean decay unchanged and corresponds to the decay control of Appendix~\ref{app:decay} (Table~\ref{tab:lm}).

\begin{table}[!ht]
\caption{MLP units in a $30$M-parameter language model after $393$M tokens (means over seeds). Random is the random-deletion tolerance of the MLP units.}
\label{tab:lm}
\begin{center}
\footnotesize
\setlength{\tabcolsep}{4pt}
\begin{tabular}{lccccc}
\toprule
Condition & Seeds & Validation loss & Nearly unused & Circuit PR & Random \\
\midrule
Adam with L2 & $5$ & $1.378$ & $0.239$ & $0.737$ & $0.422$ \\
Adam with L2, per-unit decay balance & $5$ & $1.367$ & $0.013$ & $0.988$ & $0.489$ \\
Adam with L2, ratio held & $3$ & $1.388$ & $0.307$ & $0.670$ & --- \\
AdamW, decay $0.1$ & $3$ & $1.330$ & $0.009$ & $0.99$ & $0.487$ \\
\bottomrule
\end{tabular}
\end{center}
\end{table}

Balancing each unit's decay against its gradient removes $95\%$ of the nearly unused units and lowers validation loss by $0.007$ to $0.014$ in every seed, and random-deletion tolerance rises. Equalizing the gradient while holding each unit's ratio of gradient to decay does not restore the reference but raises the unused fraction to $0.307$, possibly because the per-unit factors are noisier under minibatches. Units die outgoing column first, as in the modular-addition networks, and most of them in the first two blocks, where about half of the units end nearly unused. AdamW reaches a lower loss with fewer than $1\%$ nearly unused units, so Adam with L2 serves here to test the condition at scale and not as a recommended setting.
